% GENERATED FILE. Edit canonical lessons, then run npm run generate:books.
% canonical-source-hash: fnv1a64:2251c5020dc9a273
% canonical-lessons: TE-C245-sekarincu, TE-C245-prasaram-ceyu, TE-C245-mudrincu, TE-C245-daunlod-ceyu, TE-C245-taip-ceyu

\chapter{To collect, To broadcast, To print, To download, To type}
\label{ch:te-to-collect-to-broadcast-to-print-to-download-to-type}
\hlchaptermodality{245}

\begin{chapteropening}
\textbf{By the end of this chapter:} \emph{I can say to collect, to broadcast, to print, to download and to type.}

Complete the last lesson of chapter 245: I can say to collect, to broadcast, to print, to download and to type.
\end{chapteropening}

\section[\texorpdfstring{\te{సేకరించు} (sēkariñcu)}{సేకరించు (sēkariñcu)}]{\te{సేకరించు} (sēkariñcu) — to collect}

\label{lesson:TE-C245-sekarincu}



\begin{quote}
Before the new one: say the Telugu for to wilt, then the Telugu for to bark.
\end{quote}

\subsection*{You'll want to know: \te{సేకరించు}}
\textbf{\te{సేకరించు}} — \emph{sēkariñcu} — \textquotedblleft{}to collect\textquotedblright{}.

\begin{grammarlens}[title={What you've built}]
One more word for this chapter.
\end{grammarlens}

\subsection*{Your turn}
\begin{itemize}
\raggedright
  \item \emph{Say it:} \emph{sēkariñcu}
  \item \emph{Say it:} \emph{sēkariñcu}, once more
  \item \emph{Say it:} say \emph{sēkariñcu}
  \item \emph{Recall:} say the Telugu for to wilt, then the Telugu for to bark, then say \emph{sēkariñcu} again
\end{itemize}

\begin{tcolorbox}[breakable,colback=teal!4,colframe=teal!35!black,title={Before you move on}]
What does \te{సేకరించు} mean? (\textquotedblleft{}To collect\textquotedblright{}.) Say it once more.
\end{tcolorbox}
\section[\texorpdfstring{\te{ప్రసారం} \te{చేయు} (prasāraṁ cēyu)}{ప్రసారం చేయు (prasāraṁ cēyu)}]{\te{ప్రసారం} \te{చేయు} (prasāraṁ cēyu) — to broadcast}

\label{lesson:TE-C245-prasaram-ceyu}



\begin{quote}
Before the new one: say the Telugu for to bark, then the Telugu for to collect.
\end{quote}

\subsection*{You'll want to know: \te{ప్రసారం} \te{చేయు}}
\textbf{\te{ప్రసారం} \te{చేయు}} — \emph{prasāraṁ cēyu} — \textquotedblleft{}to broadcast\textquotedblright{}.

\begin{grammarlens}[title={What you've built}]
One more word for this chapter.
\end{grammarlens}

\subsection*{Your turn}
\begin{itemize}
\raggedright
  \item \emph{Say it:} \emph{prasāraṁ cēyu}
  \item \emph{Say it:} \emph{prasāraṁ cēyu}, once more
  \item \emph{Say it:} say \emph{prasāraṁ cēyu}
  \item \emph{Recall:} say the Telugu for to bark, then the Telugu for to collect, then say \emph{prasāraṁ cēyu} again
\end{itemize}

\begin{tcolorbox}[breakable,colback=teal!4,colframe=teal!35!black,title={Before you move on}]
What does \te{ప్రసారం} \te{చేయు} mean? (\textquotedblleft{}To broadcast\textquotedblright{}.) Say it once more.
\end{tcolorbox}
\section[\texorpdfstring{\te{ముద్రించు} (mudriñcu)}{ముద్రించు (mudriñcu)}]{\te{ముద్రించు} (mudriñcu) — to print}

\label{lesson:TE-C245-mudrincu}



\begin{quote}
Before the new one: say the Telugu for to collect, then the Telugu for to broadcast.
\end{quote}

\subsection*{You'll want to know: \te{ముద్రించు}}
\textbf{\te{ముద్రించు}} — \emph{mudriñcu} — \textquotedblleft{}to print\textquotedblright{}.

\begin{grammarlens}[title={What you've built}]
One more word for this chapter.
\end{grammarlens}

\subsection*{Your turn}
\begin{itemize}
\raggedright
  \item \emph{Say it:} \emph{mudriñcu}
  \item \emph{Say it:} \emph{mudriñcu}, once more
  \item \emph{Say it:} say \emph{mudriñcu}
  \item \emph{Recall:} say the Telugu for to collect, then the Telugu for to broadcast, then say \emph{mudriñcu} again
\end{itemize}

\begin{tcolorbox}[breakable,colback=teal!4,colframe=teal!35!black,title={Before you move on}]
What does \te{ముద్రించు} mean? (\textquotedblleft{}To print\textquotedblright{}.) Say it once more.
\end{tcolorbox}
\section[\texorpdfstring{\te{డౌన్లోడ్} \te{చేయు} (ḍaunlōḍ cēyu)}{డౌన్లోడ్ చేయు (ḍaunlōḍ cēyu)}]{\te{డౌన్లోడ్} \te{చేయు} (ḍaunlōḍ cēyu) — to download}

\label{lesson:TE-C245-daunlod-ceyu}



\begin{quote}
Before the new one: say the Telugu for to broadcast, then the Telugu for to print.
\end{quote}

\subsection*{You'll want to know: \te{డౌన్లోడ్} \te{చేయు}}
\textbf{\te{డౌన్లోడ్} \te{చేయు}} — \emph{ḍaunlōḍ cēyu} — \textquotedblleft{}to download\textquotedblright{}.

\begin{grammarlens}[title={What you've built}]
One more word for this chapter.
\end{grammarlens}

\subsection*{Your turn}
\begin{itemize}
\raggedright
  \item \emph{Say it:} \emph{ḍaunlōḍ cēyu}
  \item \emph{Say it:} \emph{ḍaunlōḍ cēyu}, once more
  \item \emph{Say it:} say \emph{ḍaunlōḍ cēyu}
  \item \emph{Recall:} say the Telugu for to broadcast, then the Telugu for to print, then say \emph{ḍaunlōḍ cēyu} again
\end{itemize}

\begin{tcolorbox}[breakable,colback=teal!4,colframe=teal!35!black,title={Before you move on}]
What does \te{డౌన్లోడ్} \te{చేయు} mean? (\textquotedblleft{}To download\textquotedblright{}.) Say it once more.
\end{tcolorbox}
\section[\texorpdfstring{\te{టైప్} \te{చేయు} (ṭaip cēyu)}{టైప్ చేయు (ṭaip cēyu)}]{\te{టైప్} \te{చేయు} (ṭaip cēyu) — to type}

\label{lesson:TE-C245-taip-ceyu}



\begin{quote}
Before the new one: say the Telugu for to print, then the Telugu for to download.
\end{quote}

\subsection*{You'll want to know: \te{టైప్} \te{చేయు}}
\textbf{\te{టైప్} \te{చేయు}} — \emph{ṭaip cēyu} — \textquotedblleft{}to type\textquotedblright{}.

\begin{grammarlens}[title={What you've built}]
One more word for this chapter.
\end{grammarlens}

\subsection*{Your turn}
\begin{itemize}
\raggedright
  \item \emph{Say it:} \emph{ṭaip cēyu}
  \item \emph{Say it:} \emph{ṭaip cēyu}, once more
  \item \emph{Say it:} say \emph{ṭaip cēyu}
  \item \emph{Recall:} say the Telugu for to print, then the Telugu for to download, then say \emph{ṭaip cēyu} again
\end{itemize}

\begin{tcolorbox}[breakable,colback=teal!4,colframe=teal!35!black,title={Before you move on}]
What does \te{టైప్} \te{చేయు} mean? (\textquotedblleft{}To type\textquotedblright{}.) Say it once more.
\end{tcolorbox}
